\section{Introducción}

Planteamiento del problema. El reto técnico de este proyecto consiste en diseñar una aplicación de línea de comandos en Python 3 que resuelva el problema clásico del cambio de monedas utilizando un enfoque algorítmico híbrido. El problema central radica en devolver una cantidad exacta de dinero empleando la menor cantidad posible de monedas a partir de un conjunto de denominaciones disponibles. Para resolver esto, el sistema debe ejecutar primero un algoritmo de ordenamiento por mezcla basado en la estrategia de divide y vencerás para organizar las monedas de mayor a menor; posteriormente, mediante una estrategia ávida implementada de forma recursiva, el programa debe seleccionar en cada llamada la denominación más grande que no exceda el monto restante, disminuyendo el saldo pendiente hasta cumplir con los casos base de paro.

\vspace{0.4cm}

Motivación. La necesidad de dar solución a este problema se basa en la importancia de estudiar y contrastar diferentes paradigmas de diseño de algoritmos para resolver tareas de optimización combinatoria en la ingeniería de software. El problema del cambio de monedas permite visualizar cómo una estrategia ávida puede ofrecer una solución óptima global de manera eficiente, siempre y cuando se trabaje bajo un sistema monetario canónico. Comprender la integración de algoritmos de ordenamiento de alta eficiencia como el ordenamiento por mezcla junto con funciones recursivas proporciona al estudiante las herramientas necesarias para analizar críticamente la complejidad computacional en tiempo y espacio, desarrollando un criterio sólido para elegir las mejores técnicas de programación ante restricciones de recursos en sistemas reales.

\vspace{0.4cm}

Objetivos. Se busca que el alumno demuestre dominio en la codificación de la estrategia de divide y vencerás para el ordenamiento descendente de las listas y en la correcta definición de los casos base dentro de la función recursiva ávida. Finalmente, se pretende realizar un análisis detallado de la complejidad temporal y espacial tanto de cada sección individual como de la solución completa, aportando los argumentos teóricos necesarios para validar el cumplimiento de los requerimientos del programa.